\section{Testing and Evaluation}\label{sec:testing}

\subsection{Automated tests}
The suite has \FactTests{} tests and uses only \code{unittest} from the standard library. It runs in
25 to 45 seconds, depending on the machine:

\begin{lstlisting}[style=text, caption={Running the tests}, label={lst:run-tests}]
$ python -m unittest discover -s tests -t . -v
...
Ran 69 tests in 25.4s
OK
\end{lstlisting}

The end-to-end tests are the ``local network simulation from several terminal windows'' of rubric
step~10, automated: \code{tests/helpers.py} starts the tracker and each peer as a separate process
with exactly the command a user would type, on free ports, and talks to them over HTTP like a
browser. The whole scenario runs twice, once with asyncio servers and once with one thread per
client. \Cref{tab:tests} summarises what is covered.

\begin{table}[H]\centering\footnotesize
\caption{Test inventory}\label{tab:tests}
\begin{tabularx}{\textwidth}{@{}l r Y@{}}
\toprule
\textbf{File} & \textbf{Tests} & \textbf{What is checked} \\
\midrule
\file{test_protocol.py} & 9 & frames split into single bytes are reassembled; oversized, empty and non-object frames are rejected; changing any AAD field or using another key fails decryption; HMAC on control packets; separate keys per purpose \\
\file{test_tracker.py} & 8 & token ownership (no take-over, no removal without token), user binding, expiry after \SI{45}{s} with a fake clock, heartbeat, salted password hashes, session cookies cannot be forged \\
\file{test_udp_files.py} & 14 & transfer without loss; 20\,\% loss on DATA \emph{and} ACK; empty file; ``\code{../../etc/}'' names; no overwrite; wrong key; receiver errors reach the sender; unreachable receiver; out-of-order and duplicate chunks; invalid chunk count; silent sender is forgotten; tampered chunk; slow start, CUBIC and RTO \\
\file{test_http.py} & 10 & query strings, pseudo-header injection, binary bodies; both server modes with 20 concurrent clients, a handler crash (500) and garbage (400); path traversal; proxy round-robin, fail-over, 404 and 502; configuration parser \\
\file{test_e2e.py} & 28 & login checked by the tracker; user, IP and ports recorded; hijack attempts refused; broadcast only to connected peers; direct messages only to the recipient; validation errors; file over lossy UDP with SHA-256 and download; crashed peer, Ctrl+C, SIGTERM, immediate restart; chat while the tracker is down; 12 peers plus 20 simultaneous registrations; sender verification, duplicates, clock skew, relay fallback \\
\bottomrule
\end{tabularx}
\end{table}

Several tests are regression tests for defects found in the original code (\cref{sec:challenges}):
the swallowed receiver error, transfers that were never forgotten, the cursor based on the sender's
clock, and the peer list that went stale.

\subsection{File transfer under packet loss}
\file{tools/bench_udp.py} sends a \SI{2}{MiB} file (\FactChunks{} chunks) between two endpoints on the
loopback interface and drops the same fraction of datagrams in \emph{both} directions, so DATA and
ACK are both lost. \Cref{tab:bench,fig:bench} give the median of three runs per loss rate.

\begin{table}[H]\centering\small
\caption{UDP file transfer of a 2\,MiB file (2468 chunks) on loopback, median of 3 runs}\label{tab:bench}
\begin{tabular}{@{}rrrrrrrc@{}}
\toprule
\textbf{Loss} & \textbf{Time (s)} & \textbf{KiB/s} & \textbf{Datagrams} & \textbf{Resent} & \textbf{Fast} & \textbf{Timeouts} & \textbf{SHA-256 ok} \\
\midrule
0\,\% & 0.41 & 4,950 & 2468 & 0 & 0 & 0 & yes \\
2\,\% & 0.52 & 3,926 & 2510 & 42 & 36 & 3 & yes \\
5\,\% & 0.85 & 2,421 & 2587 & 119 & 85 & 19 & yes \\
10\,\% & 2.65 & 774 & 2784 & 316 & 145 & 91 & yes \\
20\,\% & 11.57 & 177 & 3320 & 852 & 154 & 312 & yes \\
\bottomrule
\end{tabular}
\end{table}


\fig[fig:bench]{bench_udp.pdf}{Transfer time and number of DATA datagrams versus the loss rate}{\textwidth}

Three observations. First, the file arrived intact in every run --- the protocol is reliable at all
tested loss rates. Second, retransmission is selective: at 20\,\% loss each chunk must be sent
$1/0.8 = 1.25$ times on average just to get through once, i.e.\ about 3085 datagrams for 2468
chunks; the sender needed about 3320, only 8\,\% more than that minimum. Third, time grows much faster
than the number of datagrams: from \FactBenchZero{}\,s without loss to \FactBenchTen{}\,s at 10\,\%
and \FactBenchTwenty{}\,s at 20\,\%. The cause is visible in \cref{fig:cwnd}: without loss the window
reaches its maximum of 64 chunks in a few milliseconds of slow start, while with loss it is cut again
and again and stays between 4 and 15, and every loss that fast retransmit cannot repair costs a
timeout. This is the expected behaviour of a loss-based congestion controller --- it interprets
every loss as congestion --- and the reason why delay-based algorithms such as BBR exist.

\fig[fig:cwnd]{cwnd_trace.pdf}{Congestion window during one transfer at three loss rates}{\textwidth}

The absolute numbers are loopback numbers, measured on a two-core virtual machine: there is no real
link and the round-trip time is below a millisecond, so the limit is the Python code that encrypts,
encodes and parses every datagram (about \FactKibZero{}\,KiB/s without loss). On a real LAN the
round-trip time would make the window matter even more.

\subsection{Manual test on the LAN}
\todo{run the demo on two or three laptops in the same Wi-Fi (README section 3), note the IP
addresses, and add one screenshot of each browser; capture the traffic with Wireshark (filter
\code{tcp.port == 9101 || udp.port == 9102}) and add a screenshot showing that the chat text and file
content are not readable}
